% =============================================================================
%  2. Estudio de Mercado
% =============================================================================
\section{Estudio de Mercado}

\subsection{Definición del producto}
(Ver sección~1). En resumen, \NombreProducto{} es una solución de seguridad
automatizada para el ciclo de vida del desarrollo de software (SDLC), entregada como
suscripción en la nube y orientada al mercado centroamericano, con interfaz y soporte
en español.

\subsection{Análisis de la demanda}

\subsubsection{Distribución geográfica del mercado de consumo}
El mercado se concentra en empresas con equipos de desarrollo de software. En el
contexto regional:

\begin{itemize}
  \item \textbf{Costa Rica} es el hub tecnológico de Centroamérica: alberga más de 600
        empresas de TI y operaciones de 29 de las 100 empresas Fortune 100 (Amazon,
        Intel, IBM, HPE, P\&G, entre otras). [Ref.~8]
  \item Costa Rica y Panamá son los países centroamericanos que están desarrollando
        ecosistemas favorables para la entrada de proveedores de ciberseguridad,
        gracias a su estabilidad política y atractivo para inversión extranjera.
        [Ref.~6, Ref.~9]
  \item A nivel Latinoamérica, el mercado de ciberseguridad se concentra en Brasil
        ($\sim$39\,\% de la región), México, Argentina y Colombia; Centroamérica forma
        parte del segmento «Resto de Latinoamérica», de menor tamaño pero con alto
        potencial de crecimiento. [Ref.~7, Ref.~10]
\end{itemize}

\noindent\textbf{Segmentación del mercado meta de \NombreProducto:}
\begin{enumerate}
  \item Empresas de desarrollo de software y \emph{software factories} (Costa Rica
        concentra la mayoría).
  \item Sector financiero (BFSI): bancos, fintech y cooperativas, el vertical de mayor
        inversión en ciberseguridad de la región. [Ref.~7, Ref.~10]
  \item Centros de servicios compartidos y operaciones nearshore de multinacionales.
  \item Sector gobierno, impulsado tras los ataques de ransomware de 2022. [Ref.~9]
\end{enumerate}

\subsubsection{Comportamiento histórico de la demanda}
La demanda de soluciones DevSecOps ha mostrado crecimiento exponencial sostenido a
nivel global:

\begin{table}[ht]
  \centering
  \begin{tabular}{cc}
    \toprule
    \textbf{Año} & \textbf{Tamaño mercado global DevSecOps (USD miles de millones)} \\
    \midrule
    2024 & $\sim$8.8 -- 9.7 \\
    2025 & $\sim$9.0 -- 10.3 \\
    2026 & $\sim$10.9 -- 11.7 \\
    \bottomrule
  \end{tabular}
  \caption{Fuentes: [Ref.~1, Ref.~2, Ref.~3, Ref.~5]}
\end{table}

Factores que han impulsado la demanda histórica: aumento de ciberataques y filtraciones
de datos, adopción creciente de metodología DevOps, mayores exigencias de cumplimiento
regulatorio y la popularización de servicios en la nube. [Ref.~4]

En Centroamérica, el detonante más fuerte fue el \textbf{ataque de ransomware
Conti/Hive a Costa Rica en 2022}, que provocó un estado de emergencia nacional,
paralizó sistemas de Hacienda y de la CCSS, y generó pérdidas estimadas en USD 30
millones por día. Este evento elevó drásticamente la conciencia y el gasto en
ciberseguridad en la región. [Ref.~9]

\subsubsection{Proyección de la demanda}
Las proyecciones de las principales firmas de investigación coinciden en un crecimiento
de dos dígitos:

\begin{table}[ht]
  \centering
  \begin{tabularx}{\linewidth}{Yl}
    \toprule
    \textbf{Indicador} & \textbf{Valor} \\
    \midrule
    CAGR global DevSecOps (2026--2031/2035) & 22\,\% -- 28\,\% (conservadoras: 13\,\% -- 14\,\%) \\
    Mercado global DevSecOps proyectado 2030--2031 & USD 27 -- 30 mil millones \\
    Segmento de mayor crecimiento & PYMEs (SMEs) y despliegues en la nube \\
    CAGR del segmento Servicios & $\sim$25\,\% -- 27\,\% \\
    \bottomrule
  \end{tabularx}
  \caption{Fuentes: [Ref.~1, Ref.~4, Ref.~5]}
\end{table}

Para el mercado latinoamericano de ciberseguridad, la proyección es crecer de
USD 11.7 mil millones (2025) a USD 18.4 mil millones (2030), con un CAGR de 9.5\,\%.
El segmento de seguridad en la nube y los servicios gestionados son los de mayor
crecimiento. [Ref.~7, Ref.~10]

\noindent\textbf{Proyección específica para \NombreProducto{} (escenario base
conservador, captación gradual y modesta):}

\begin{table}[ht]
  \centering
  \begin{tabular}{cccc}
    \toprule
    \textbf{Año} & \textbf{Clientes (empresas)} & \textbf{Devs prom./cliente} & \textbf{Suscripciones (seats)} \\
    \midrule
    1 & 4  & 10 & 40  \\
    2 & 9  & 11 & 99  \\
    3 & 16 & 12 & 192 \\
    4 & 25 & 13 & 325 \\
    5 & 38 & 14 & 532 \\
    \bottomrule
  \end{tabular}
\end{table}

\emph{Estas cifras suponen una startup pequeña y bootstrapped, que crece de boca en
boca y con pocos recursos comerciales — un escenario prudente y defendible para un
mercado emergente como el centroamericano.}

\subsubsection{Tabulación de datos de fuentes primarias (propuesta metodológica)}
Como parte del estudio de mercado se plantea la aplicación de una \textbf{encuesta a
fuentes primarias} dirigida al mercado meta (empresas con equipos de desarrollo de
software en Costa Rica: software factories, áreas de TI de bancos y fintech, startups
tecnológicas y centros de servicios de multinacionales). A continuación se describe el
instrumento propuesto y se presenta una \textbf{tabulación con valores ilustrativos}
que ejemplifica el tipo de resultados esperados.

\noindent\textbf{Ficha técnica propuesta del estudio:}
\begin{itemize}
  \item \emph{Instrumento:} cuestionario estructurado en línea (Google Forms / Microsoft Forms).
  \item \emph{Población objetivo:} empresas con equipos de desarrollo de software en Costa Rica.
  \item \emph{Muestra sugerida:} 30--40 empresas del sector TI.
  \item \emph{Método de contacto:} correo, LinkedIn y comunidades de desarrolladores.
  \item \emph{Tiempo de respuesta estimado:} 3 minutos por participante.
\end{itemize}

\noindent\textbf{Modelo de tabulación (valores ilustrativos a modo de ejemplo):}

\begin{table}[ht]
  \centering
  \begin{tabularx}{\linewidth}{Y Y c}
    \toprule
    \textbf{Pregunta} & \textbf{Respuesta} & \textbf{\%} \\
    \midrule
    ¿Su equipo usa pipelines CI/CD?                       & Sí                        & 78\,\% \\
    ¿Realiza pruebas de seguridad automatizadas hoy?      & No / parcialmente         & 65\,\% \\
    ¿Ha sufrido un incidente de seguridad (últimos 2 años)? & Sí                      & 42\,\% \\
    ¿Pagaría por una solución SaaS de seguridad integrada? & Sí                       & 58\,\% \\
    Rango de pago mensual aceptable por desarrollador     & USD 15 -- 40              & 61\,\% \\
    Principal barrera de adopción                         & Presupuesto / falta de talento & 70\,\% \\
    \bottomrule
  \end{tabularx}
\end{table}

\emph{Nota: los porcentajes anteriores son valores de referencia ilustrativos que
muestran el tipo de información que el instrumento permitiría obtener; al aplicar la
encuesta en campo, estos datos se sustituirían por los resultados reales tabulados.}

\noindent\textbf{Lectura de los datos ilustrativos:} el patrón esperado es consistente
con el análisis de fuentes secundarias: la mayoría de los equipos ya usa CI/CD, pero
pocos realizan pruebas de seguridad automatizadas, y la principal barrera es el
presupuesto y la falta de talento especializado.

\noindent\textbf{Conclusión de la demanda:} existe una demanda real, creciente e
insatisfecha, especialmente en PYMEs que tienen pipelines CI/CD pero carecen de
especialistas en seguridad — exactamente el perfil que un modelo SaaS asequible y en
español puede capturar.

\subsection{Análisis de la oferta}

\subsubsection{Características de los principales productores / competencia}
Los competidores globales son consolidados y dominan el mercado mundial:

\begin{table}[ht]
  \centering
  \begin{tabularx}{\linewidth}{l Y Y}
    \toprule
    \textbf{Competidor} & \textbf{Fortaleza} & \textbf{Debilidad frente a \NombreProducto} \\
    \midrule
    \textbf{Snyk} & Líder en SCA y experiencia de desarrollador & Precio alto, soporte y facturación en inglés/USD \\
    \textbf{Checkmarx} & SAST robusto, enfoque enterprise & Complejo y costoso para PYMEs \\
    \textbf{GitLab (Ultimate)} & Seguridad integrada en su plataforma & Solo sirve a quienes usan GitLab \\
    \textbf{Veracode} & Cumplimiento y reportería & Orientado a grandes corporaciones \\
    \textbf{Palo Alto (Prisma Cloud)} & Suite completa cloud & Sobredimensionado y caro para el mercado local \\
    \bottomrule
  \end{tabularx}
  \caption{Fuentes: [Ref.~1, Ref.~4]}
\end{table}

A nivel regional, han surgido firmas latinoamericanas relevantes como \textbf{Hackmetrix,
Conviso, UNXPOSE, Clavis Security y Cipher}, que compiten agregando servicios y
conocimiento local. [Ref.~10]

\noindent\textbf{Ventaja competitiva de \NombreProducto:}
\begin{itemize}
  \item Precio accesible y facturación en moneda local/USD adaptada a PYMEs.
  \item Interfaz, documentación y soporte técnico en español, en zona horaria local.
  \item Plantillas de cumplimiento mapeadas a la Ley 8968 de Costa Rica y normativa regional.
  \item Onboarding asistido para empresas sin especialistas en seguridad.
\end{itemize}

\subsubsection{Proyección de la oferta}
La oferta global de soluciones DevSecOps se expande de forma acelerada, con las
soluciones tipo plataforma representando $\sim$70\,\% del mercado. [Ref.~1] Sin embargo,
\textbf{la oferta especializada y localizada para PYMEs centroamericanas es escasa}:
los grandes proveedores priorizan grandes cuentas en inglés y los proveedores locales
se enfocan más en servicios de consultoría que en producto SaaS propio. Esto deja un
nicho desatendido que \NombreProducto{} puede ocupar en los próximos 3--5 años antes
de que la competencia internacional regionalice su oferta.

\subsection{Importaciones del producto o servicio}
Al ser un SaaS, no hay «importación» física, pero sí una dependencia tecnológica
importante: prácticamente toda la oferta consumida actualmente en Costa Rica y
Centroamérica es \textbf{importada} (Snyk, Checkmarx, GitHub Advanced Security, etc.),
facturada en dólares desde el exterior. Esto implica:

\begin{itemize}
  \item Salida de divisas por suscripciones a proveedores extranjeros.
  \item Costos de infraestructura cloud (AWS/Azure/GCP) que también se pagan en USD.
  \item \NombreProducto{} usaría infraestructura de hyperscalers (importada), pero
        generaría valor agregado, empleo y propiedad intelectual local, sustituyendo
        parcialmente la importación del servicio final.
\end{itemize}

\subsection{Análisis de precios}

\subsubsection{Determinación del costo promedio}
El precio de mercado de soluciones DevSecOps se cobra típicamente \emph{por
desarrollador/mes}. Rango de referencia internacional: USD 25 -- 60 por
desarrollador/mes en planes profesionales.

\noindent\textbf{Estructura de precios propuesta para \NombreProducto:}

\begin{table}[ht]
  \centering
  \begin{tabularx}{\linewidth}{l c Y}
    \toprule
    \textbf{Plan} & \textbf{Precio (USD/dev/mes)} & \textbf{Público} \\
    \midrule
    Starter    & 18 & PYMEs, startups (hasta 10 devs) \\
    Business   & 30 & Empresas medianas \\
    Enterprise & 45 + soporte dedicado & Banca, gobierno, multinacionales \\
    \bottomrule
  \end{tabularx}
\end{table}

\noindent\textbf{Precio promedio ponderado estimado:} $\sim$USD 28 por desarrollador/mes.

\subsubsection{Análisis histórico y proyección de precios}
Históricamente, los precios por \emph{seat} han tendido a mantenerse o subir levemente,
pero la aparición de modelos de \textbf{precio por consumo} (pay-as-you-go) está
reduciendo la barrera de entrada, sobre todo para PYMEs. [Ref.~5] La volatilidad
cambiaria en la región empuja a los clientes hacia suscripciones en USD y planes
flexibles en lugar de licencias plurianuales. [Ref.~11]

\noindent\textbf{Proyección de precio promedio de \NombreProducto:} mantener USD 28/seat
los primeros 2 años (penetración), con incrementos anuales de $\sim$3--5\,\% a partir
del año 3, alineados a la inflación y a mejoras de producto (IA de remediación).

\subsection{Canales de comercialización y distribución}

\subsubsection{Descripción de los canales de distribución}
Al ser un producto digital (SaaS), la distribución es 100\,\% en línea:

\begin{enumerate}
  \item \textbf{Venta directa (self-service):} registro y prueba gratuita desde el
        sitio web; el cliente conecta su repositorio y empieza a usar la plataforma.
  \item \textbf{Venta consultiva B2B:} equipo comercial para cuentas medianas y grandes
        (banca, gobierno), con demos personalizadas.
  \item \textbf{Marketplaces de nube:} publicación en AWS Marketplace y Microsoft Azure
        Marketplace para llegar a clientes que ya compran ahí.
  \item \textbf{Alianzas con partners:} integradores de TI y consultoras locales que
        revenden o recomiendan la plataforma (modelo de canal con comisión).
  \item \textbf{Marketing de contenido y comunidad:} blog técnico, webinars y presencia
        en comunidades de desarrolladores para generar \emph{leads}.
\end{enumerate}
